% Options for packages loaded elsewhere
\PassOptionsToPackage{unicode}{hyperref}
\PassOptionsToPackage{hyphens}{url}
\documentclass[
  ignorenonframetext,
]{beamer}
\newif\ifbibliography
\usepackage{pgfpages}
\setbeamertemplate{caption}[numbered]
\setbeamertemplate{caption label separator}{: }
\setbeamercolor{caption name}{fg=normal text.fg}
\beamertemplatenavigationsymbolsempty
% remove section numbering
\setbeamertemplate{part page}{
  \centering
  \begin{beamercolorbox}[sep=16pt,center]{part title}
    \usebeamerfont{part title}\insertpart\par
  \end{beamercolorbox}
}
\setbeamertemplate{section page}{
  \centering
  \begin{beamercolorbox}[sep=12pt,center]{section title}
    \usebeamerfont{section title}\insertsection\par
  \end{beamercolorbox}
}
\setbeamertemplate{subsection page}{
  \centering
  \begin{beamercolorbox}[sep=8pt,center]{subsection title}
    \usebeamerfont{subsection title}\insertsubsection\par
  \end{beamercolorbox}
}
% Prevent slide breaks in the middle of a paragraph
\widowpenalties 1 10000
\raggedbottom
\AtBeginPart{
  \frame{\partpage}
}
\AtBeginSection{
  \ifbibliography
  \else
    \frame{\sectionpage}
  \fi
}
\AtBeginSubsection{
  \frame{\subsectionpage}
}
\usepackage{iftex}
\ifPDFTeX
  \usepackage[T1]{fontenc}
  \usepackage[utf8]{inputenc}
  \usepackage{textcomp} % provide euro and other symbols
\else % if luatex or xetex
  \usepackage{unicode-math} % this also loads fontspec
  \defaultfontfeatures{Scale=MatchLowercase}
  \defaultfontfeatures[\rmfamily]{Ligatures=TeX,Scale=1}
\fi
\usepackage{lmodern}
\ifPDFTeX\else
  % xetex/luatex font selection
\fi
% Use upquote if available, for straight quotes in verbatim environments
\IfFileExists{upquote.sty}{\usepackage{upquote}}{}
\IfFileExists{microtype.sty}{% use microtype if available
  \usepackage[]{microtype}
  \UseMicrotypeSet[protrusion]{basicmath} % disable protrusion for tt fonts
}{}
\makeatletter
\@ifundefined{KOMAClassName}{% if non-KOMA class
  \IfFileExists{parskip.sty}{%
    \usepackage{parskip}
  }{% else
    \setlength{\parindent}{0pt}
    \setlength{\parskip}{6pt plus 2pt minus 1pt}}
}{% if KOMA class
  \KOMAoptions{parskip=half}}
\makeatother
\setlength{\emergencystretch}{3em} % prevent overfull lines
\providecommand{\tightlist}{%
  \setlength{\itemsep}{0pt}\setlength{\parskip}{0pt}}
\usepackage{bookmark}
\IfFileExists{xurl.sty}{\usepackage{xurl}}{} % add URL line breaks if available
\urlstyle{same}
\hypersetup{
  hidelinks,
  pdfcreator={LaTeX via pandoc}}

\author{\texorpdfstring{}{}}
\date{}

\begin{document}

\begin{frame}{TQNNs and ZX-Calculus}
\protect\phantomsection\label{sec:quantum-active-inference}
The preceding sections have established that categorical string
diagrams---from DisCoCat's pregroup derivations
(\autoref{sec:categorical-semantics}) through enriched hom-values
(\autoref{sec:enriched-categories}) to topos-theoretic transfer
(\autoref{sec:topos-theory})---provide a unified diagrammatic language
for case-theoretic reasoning. This section extends the framework into
its natural quantum generalization: topological quantum neural networks
(TQNNs), the ZX-calculus, and sheaf-theoretic quantum semantic
communication. The central observation is that the same
monoidal-categorical architecture that underlies DisCoCat string
diagrams also underlies quantum circuits, TQFT cobordisms, and quantum
information flow---making active inference on case-marked relational
structure a quantum topological computation.

\begin{block}{QNNs as Spin-Networks}
\protect\phantomsection\label{qnns-as-spin-networks}
\begin{block}{TQFT as the Forward Pass: Reshetikhin--Turaev Invariants
Compute Network Amplitudes}
\protect\phantomsection\label{tqft-as-the-forward-pass-reshetikhinturaev-invariants-compute-network-amplitudes}
Marcianò, Fields, and Glazebrook show that quantum neural networks
(QNNs) admit a topological reformulation using spin-networks
{[}@fields2022tqnn{]}. Any QNN layer can be represented as a graph whose
edges carry representation labels (spins) and whose vertices carry
intertwiners---precisely the data defining a spin-network in a
3-dimensional topological quantum field theory (TQFT):

\begin{quote}
``Quantum Neural Networks (QNNs) can be mapped onto spin-networks, with
the consequence that the level of analysis of their operation can be
carried out on the side of Topological Quantum Field Theory (TQFT).''
{[}@fields2022tqnn{]}
\end{quote}

This reformulation has three structural consequences. First, the network
becomes a topological diagram---spin-network or ribbon graph---evaluated
by a continuous TQFT functor; edges encode information flow and nodes
encode transformation. Second, the TQFT evaluation assigns boundary
Hilbert spaces to the diagram via the Reshetikhin--Turaev and
Turaev--Viro invariants, playing the role of the neural forward pass:
quantum amplitudes propagate through the topological structure. Third,
information flow is encoded in the topology of the wiring rather than in
any fixed geometric embedding, giving the architecture inherent
robustness to continuous deformation {[}@fields2023tensor{]}.
\end{block}

\begin{block}{TQNNs Are Universal}
\protect\phantomsection\label{tqnns-are-universal}
Fields and collaborators further demonstrate that TQNNs are universal
quantum computers by identifying the Reshetikhin--Turaev invariant of a
TQNN with a Turaev--Viro quantum error-correcting code:

\begin{quote}
``TQNNs enable universal quantum computation, using the
Reshetikhin-Turaev and Turaev-Viro models to show how TQNNs implement
quantum error-correcting codes.'' {[}@fields2025amplituhedra{]}
\end{quote}

The universality result is established via the concept of an
\emph{execution trace} for a quantum computation, leading to the
representation of TQNNs in terms of the positive geometries provided by
amplituhedra---a deep connection between quantum computation, scattering
amplitudes, and topological combinatorics.
\end{block}

\begin{block}{QRFs Select the Measurement Basis}
\protect\phantomsection\label{qrfs-select-the-measurement-basis}
Fields and Glazebrook's work on quantum reference frames (QRFs) and
holographic screens provides additional algebraic structure
{[}@fields2021qrf{]}. A holographic screen---the information boundary
between two interacting quantum systems---carries a qubit array encoding
their interaction. The key insight is that QRFs deployed to identify
systems and select pointer states induce decoherence, breaking the
symmetry of the holographic encoding in an observer-relative way. This
symmetry-breaking is precisely the mechanism by which a TQNN
``observes'' its input: the choice of QRF determines the basis in which
the spin-network is evaluated.

For case-theoretic reasoning, this connects to the grammatical observer
problem: a parser or comprehender selecting a case-assignment frame for
a sentence is analogous to deploying a QRF that fixes the pointer basis
for a quantum measurement on a holographic screen.
\end{block}
\end{block}

\begin{block}{ZX-Calculus: Topological String Diagrams Where Graph
Rewrites Are Quantum Proofs}
\protect\phantomsection\label{zx-calculus-topological-string-diagrams-where-graph-rewrites-are-quantum-proofs}
\begin{block}{String Diagrams for Quantum Processes}
\protect\phantomsection\label{string-diagrams-for-quantum-processes}
The powerful ZX-calculus provides a diagrammatic language for quantum
circuits, representing them as string diagrams in a symmetric monoidal
category of finite-dimensional Hilbert spaces and linear maps
{[}@kissinger2020zx{]}:

\begin{quote}
``The ZX-calculus is a graphical language for reasoning about quantum
computations and circuits\ldots{} it can represent any linear map, and
can be considered a diagrammatically complete generalization of the
usual circuit representation.'' {[}@coeckekissinger2017{]}
\end{quote}

Three structural features connect ZX to case-theoretic diagrams:

\begin{itemize}
\tightlist
\item
  \textbf{Diagrams as morphisms}: ZX-diagrams are string diagrams in a
  \(\dagger\)-compact closed category. Wires represent objects (qubits),
  spiders and boxes represent morphisms, and composition/tensor product
  correspond to vertical/horizontal concatenation. Only topology
  matters, not geometry.
\item
  \textbf{Deterministic circuit extraction via generalized flow}:
  Kissinger and van de Wetering show that quantum circuits map to
  ZX-diagrams, undergo graph-theoretic rewriting, and extract as
  optimized circuits---with topological abstraction preserving the
  invariants needed for optimization {[}@kissinger2020zx{]}.
\item
  \textbf{Category-theoretic semantics}: A ZX-diagram's semantics are
  determined entirely by how components are wired---the same
  compositional principle underlying DisCoCat and the case categories of
  \autoref{sec:case-systems}.
\end{itemize}
\end{block}

\begin{block}{Cups Are Spiders}
\protect\phantomsection\label{cups-are-spiders}
The structural parallel between pregroup grammar diagrams and
ZX-diagrams is not accidental. Both are instances of the same
mathematical object: morphisms in a compact closed monoidal category
with functorial semantics into Hilbert spaces. In DisCoCat, the functor
assigns to each grammatical type a vector space and to each derivation a
linear map computing sentence meaning {[}@coecke2010mathematical{]}. In
ZX, the standard semantics functor assigns to each spider/Hadamard
configuration a linear map in \textbf{FHilb}. The shared categorical
architecture means:

\begin{enumerate}
\tightlist
\item
  \textbf{Pregroup contractions (cups) and ZX spiders} are both
  instances of the same algebraic operation: evaluation morphisms in a
  compact closed category.
\item
  \textbf{DisCoCat normal forms} and \textbf{ZX simplifications} are
  both applications of the same rewriting theory: equational reasoning
  modulo the axioms of a compact closed category.
\item
  \textbf{The snake equation} (Cap \(\circ\) Cup = identity) that
  grounds all pregroup type reductions
  (\autoref{sec:categorical-semantics}) is a special case of the spider
  fusion rule in ZX.
\end{enumerate}

This means that case-theoretic DisCoCat derivations can, in principle,
be compiled into ZX circuits and executed on quantum hardware---a
connection already exploited by the lambeq quantum NLP pipeline
{[}@lorenz2023lambeq{]}.
\end{block}
\end{block}

\begin{block}{One Diagram, Three Interpretations}
\protect\phantomsection\label{one-diagram-three-interpretations}
\begin{block}{Common Language: Ribbon and Tensor Diagrams}
\protect\phantomsection\label{common-language-ribbon-and-tensor-diagrams}
Both ZX-diagrams and TQNNs are topological string diagrams evaluated by
monoidal functors into the category of Hilbert spaces and linear maps.
The alignment becomes explicit when stated precisely:

\begin{itemize}
\tightlist
\item
  \textbf{For TQNNs}: A 3-dimensional TQFT functor from a cobordism or
  skein category to \textbf{Hilb} assigns to each spin-network/ribbon
  graph a linear map representing the TQNN computation. The underlying
  topological skein theory treats network layers as \emph{ribbon graphs}
  whose evaluation via Reshetikhin--Turaev and Turaev--Viro invariants
  gives quantum processes implementing computation and error correction
  {[}@fields2022tqnn; @fields2025amplituhedra{]}.
\item
  \textbf{For ZX}: The standard semantics functor from the free
  \(\dagger\)-compact category generated by Z/X spiders, Hadamards, etc.
  into \textbf{FHilb} assigns each ZX-diagram a linear map
  {[}@kissinger2020zx{]}.
\item
  \textbf{For DisCoCat}: The meaning functor from the pregroup grammar
  category to \textbf{FdVect} assigns to each grammatical derivation a
  multilinear map computing compositional meaning
  {[}@coecke2010mathematical{]}.
\end{itemize}

Up to choice of labels and normalization, all three are graphical
calculi for monoidal categories whose morphisms are quantum (or
quantum-like) processes. A layer of a topological quantum flow network
can be modeled as a ZX-diagram fragment whose input and output boundary
wires are the ``feature spaces'' (Hilbert spaces) at successive
processing steps, while internal spiders encode unitary/non-unitary
channels that realize synaptic transformations.
\end{block}

\begin{block}{Generalized Flow Guarantees Causal Order}
\protect\phantomsection\label{generalized-flow-guarantees-causal-order}
The \emph{generalized flow} condition used to guarantee deterministic
circuit extraction from ZX-diagrams is a graph-theoretic constraint that
ensures a well-defined causal ordering of operations
{[}@kissinger2020zx{]}. This mirrors the requirement in TQNNs that the
diagram encode a consistent \emph{execution trace} of a quantum
computation. Fields and colleagues make this connection explicit in
their TQNN--amplituhedron correspondence:

\begin{quote}
``\ldots this formal correspondence is stated by Theorem 2 whose proof
draws upon the concept of execution trace for a quantum
computation\ldots{} and thus leads to representing a TQNN in terms of
the positive geometries as provided by amplituhedra.''
{[}@fields2025amplituhedra{]}
\end{quote}

A topological quantum flow neural network can therefore be regarded as a
ZX-style circuit where the graphical calculus is enriched to a 3D TQFT
skein theory, but the \emph{abstract type} of object---a topological
string diagram with functorial semantics---remains the same.
\end{block}
\end{block}
\end{frame}

\end{document}
